\chapter{Wissen abbilden mittels Ontologien}
\label{chap:ontologien}

Die bisherigen Kapitel bieten einen groben Überblick über Objekte und deren Anwendung. Um bei der Wissensmodellierung von Nutzen zu sein, müssen sie aber in eine geeignete Form gebracht werden.\\
Wissen wird im knowledge engineering meist in Form von Ontologien abgelegt.

Der Begriff Ontologie wird in verschiedenen wissenschaftlichen Bereichen verwendet, so zum Beispiel in der Philosophie, der Psychologie und der Informatik. In diesem Dokument soll nur der Aspekt der Informatik erklärt werden.

Das folgende Kapitel basiert auf dem Artikel ``Was bedeutet eigentlich Ontologie?''~\cite{IspekOntoBedeutung} sowie dem Artikel ``Eine kurze Geschichte der Ontologie''~\cite{ISpekOntoGeschichte}.

In der Informatik steht Ontologie für ``\ldots eine formale Beschreibung des Wissens in einer Domäne in der Form von Konzepten der Domäne, deren Beziehung untereinander und der Eigenschaft dieser Konzepte und Beziehungen, sowie der in der Domäne gültigen Axiome und Prinzipien.''~\cite[S.310]{ISpekOntoGeschichte}.

Zum besseren Verständnis wird diese Definition im folgenden Abschnitt kurz untersucht:
\begin{itemize}
    \item \textit{Domäne}

        Unter einer Domäne versteht man einen Ausschnitt der Welt, dessen Grenzen klar definiert sind. Die Domäne und ihre Grenzen werden durch den Anwendungsfall festgelegt.

    \item \textit{Konzepte}

        Konzepte werden in anderen Teilen dieses Dokumentes als Objekte oder Klassen bezeichnet. Dabei kann es sich um materielle Konzepte, wie zum Beispiel Teile eines Wagens oder um immaterielle Konzepte, wie zum Beispiel die Lösungssuche, handeln.

    \item \textit{Beziehungen}

        Objekte stehen in Beziehungen zueinander. Es werden die Beziehungen ``ist ein'' und ``Instanz von'' unterschieden. In unserem Beispiel sind eine Abenteuerreise eine Reise und der Seilpark Balmberg eine Instanz eines Ausfluges ($Abenteuerreise\;istEin\;Ausflug$ und $SeilparkBalmberg\;instanzVon\;Ausflug$).\\
        Die Eigenschaften eines Objektes werden auch als Beziehungen dargestellt (eine Reise hat einen Reiseführer).

    \item \textit{Axiome und Prinzipien}

        Damit werden die in der Domäne vorhandenen Regeln bezeichnet.
\end{itemize}

Mit diesen Elementen sind alle wichtigen Punkte einer Ontologie abgedeckt. Dabei ist es, wie in der Definition festgelegt, wichtig, dass die gesamte Abbildung in einer formalen Sprache beschrieben wird.

\subsection{Anwendung von Ontologien}
\label{subsec:ontologien_onto_anwendung}
Ontologien werden für wissensbasierte Anwendungen verwendet. Konkret werden sie hauptsächlich dort verwendet, wo Semantik zur Formulierung von Informationen genutzt wird. Eines der bekanntesten Beispiele ist das semantische Web. Im Gegensatz zum syntaktischen Web versucht das semantische Web Zeichen nicht nur abzugleichen, sondern auch ein Verständnis und Schlussfolgerungen einzubauen.

Die Aufgabe von Ontologien ist die Ermöglichung und Verbesserung von Kommunikation zwischen Computeranwendungen untereinander und zwischen Computeranwendungen und Mensch. Wichtig: Es darf sich immer nur um eine bestimmte Wissensdomäne (Domain Ontology) handeln.\\
Einige Berufsgruppen haben bereits begonnen, ihr Wissen mit Hilfe von standardisierten Sprachen abzubilden. Diese Sprachen, wie die in Kapitel~\ref{chap:owl} vorgestellte Sprache OWL, wurden vom W3C-Konsortium (World Wide Web Consortium) standardisiert.

\subsection{Anwendung des semantischen Webs}
\label{subsec:ontologien_onto_SemantikWebAnwendung}
Eine Anwendung für das semantische Web besteht aus der Gewinnung von zusätzlichem Wissen, was anhand einer Wissensbasis in Form einer Wissensdatenbank und einer Inferenzmaschine geschieht. Ein Ontologieschema legt in der Wissensbasis fest, welche Arten von Aussagen möglich sind. Die Aussagen enthalten das konkrete Wissen. Sowohl das Schema, wie auch die Wissensabbildung werden in einer formalen Sprache, wie z.B. OWL oder RDFS, abgebildet.

\begin{lstlisting}[caption={Beispiel einer Aussage in einer Wissensbasis.}]
Eine Abenteuerreise ist eine Reise.
\end{lstlisting}

In der einfachsten Anwendung kann das abgelegte Wissen direkt abgefragt werden. Bei komplexeren Anfragen sind die Funktionalitäten einer \textit{Inferenzmaschine} notwendig. Mittels Inferenzregeln ist diese fähig Schlussfolgerungen zu ziehen und dadurch neue, komplexe Aussagen zu treffen.

Die Transitivitätsregel ist eine einfache, von der Inferenzmaschine verwendete Regel: ``Wird eine Relation $p$ als transitiv deklariert und es gilt $x p y$ und $y p z$, dann kann $x p z$ gefolgert werden.''~\cite[Seite 289]{IspekOntoBedeutung}
Dem Entwickler ist es möglich, neben den vorhandenen Regeln, wissensdomänenspezifische Regeln zu spezifizieren. Die Inferenzmaschine berücksichtigt diese bei ihrer Auswertung.

\begin{lstlisting}[caption={Beispiel einer Regel in einer Wissensbasis.}]
Handelt es sich bei einem Objekt um eine Reise, benötigt das Objekt mindestens vier Teilnehmer/innen und ist das Objekt teambildend, so handelt es sich bei dem Objekt auch um einen Teamevent.
\end{lstlisting}

Damit eine Inferenzmaschine Schlüsse ziehen kann, muss die Wissensbasis in einer formalen Sprache vorliegen. Sie darf keine syntaktischen Fehler enthalten. Das Schema und die Aussagen werden von der Maschine geladen und intern als Graph gespeichert (siehe Kapitel~\ref{chap:wissensrepFormen}).\\
Durch einen Algorithmus zum Abgleich von Graphen werden Abfragen implementiert. Durch wiederholtes Anwenden von Regeln können \textit{Schlussfolgerungen} gezogen werden.\\ Die Anwendung dieser Regeln erfolgt wiederum mit Hilfe von Algorithmen zum Abgleich von Graphen. Die dadurch entstandenen Aussagen werden der Wissensbasis hinzugefügt und stehen für Abfragen zur Verfügung.

\noindent\rule[1ex]{\textwidth}{1pt}
\begin{wrapfigure}[8]{l}{0.1\textwidth}
    \vspace{-14pt}
    \includegraphics[width=0.1\textwidth]{bilder/owl.png}
\end{wrapfigure}
Wir brauchen Ontologien. Dabei spielt die Wahl der Domäne eine grosse Rolle. Bei unserer Wissensmodellierung haben wir gelernt, das sich nicht jeder Themenbereich als Domäne eignet.
Bei Umsetzung der ursprünglichen Problemdomäne, nämlich dem Erlernen des Programmierens anhand der Programmiersprache Prolog, mussten wir feststellen, dass diese nicht als Domäne geeignet ist. Zum Beispiel lässt sich das Konzept der Rekursion in einer Ontologie zwar lexikal abbilden, der Mehrwert (einer Ontolgie gegenüber der herkömmlichen Wissensabbildung) in Form von Inferenz, ist jedoch nicht gegeben. Die Möglichkeiten einer Wissensmodellierung anhand einer Ontologie sind bei Themen, die Inferenz und damit Schlussfolgerungen zulassen, viel grösser.
\noindent\rule[1ex]{\textwidth}{1pt}

Im nachfolgenden Kapitel wird eine Erklärung gegeben, wie Inferenz und Resolution zum Ziehen von Schlüssen verwendet werden.
